\documentclass[letterpaper,11pt]{article}
\usepackage[top=0.8in, bottom=0.8in, left=0.9in, right=0.9in]{geometry}
\usepackage[hidelinks]{hyperref}
\usepackage{lmodern}
\usepackage{parskip}
\usepackage{microtype}
\setlength{\parskip}{6pt}
\setlength{\parindent}{0pt}
\begin{document}
\begin{flushleft}
\textbf{\large Ajay Venkatesh} \\
Brooklyn, NY \\
+1 516 585 9013 \\
\href{mailto:av3855@nyu.edu}{av3855@nyu.edu}
\end{flushleft}
\vspace{10pt}
May 21, 2026
\vspace{6pt}

Hiring Manager \\
Lam Research

\vspace{10pt}

Dear Hiring Manager,

My name is Ajay Venkatesh. I'm finishing my M.S. in Computer Engineering at NYU, with production software experience at PwC in Bengaluru, a summer SDE internship at Amazon, and ongoing on-campus work at NYU's Novel AI Technologies. I'm writing about the MES/IoT Developer role at the Tualatin R\&D Center.

The role's combination, automation software for equipment interfaces, communication frameworks, and IoT-enabled solutions for R\&D labs, sits where my engineering work and graduate coursework intersect. My coursework spans \textbf{Computer Architecture}, \textbf{System Design}, \textbf{Data Structures}, and \textbf{Big Data} alongside \textbf{Advanced Java}, exactly the systems-and-architecture foundation that semiconductor automation work needs.

On building reliable backend systems for connected devices: at NYU's Novel AI Technologies I architected an AI-powered health platform for iOS and Android that ingests real-time wearable telemetry from connected devices into cloud-hosted processing, exactly the kind of device-to-cloud data flow MES/IoT work demands. At Amazon I built a high-throughput, event-driven processing system on \textbf{AWS ECS Fargate, SQS, SNS} that decoupled distributed message handling under variable load, with provisioning through \textbf{AWS CDK} (Lambda, API Gateway, DynamoDB, IAM, CloudWatch) and multi-stage CI/CD.

On Java services and multithreading: at Amazon I shipped backend services in \textbf{Java} (Coral, Smithy, Dagger, CBOR) with low-level design patterns and unit / integration testing through \textbf{Mockito}. I've worked on concurrent transaction handling in Python projects through threading and synchronized blocks. The asynchronous-processing discipline carries over directly to MES/IoT communication frameworks.

Stack-wise: \textbf{Java}, \textbf{Python}, \textbf{C++}, \textbf{SQL}, \textbf{REST APIs}, \textbf{Docker}, \textbf{Kubernetes}, plus design patterns (OOP, MVC, DAO) and multithreaded programming from coursework and projects. Comfortable across Linux environments and at home with the integration-style work that connects equipment to higher-level systems.

What pulls me toward Lam Research is the layered engineering. Semiconductor R\&D automation runs at the seam between hardware and software where small reliability improvements compound across labs and tools, and contributing to innovation that feeds back into manufacturing is engineering with real downstream impact.

I'd welcome the chance to discuss further. Thanks for considering my application.

\vspace{12pt}
Sincerely, \\
Ajay Venkatesh

\end{document}